Positive row scaling preserves a mixed-integer program in exact arithmetic, but changes its
finite-precision representation and the original-unit meaning of residual tolerances.
We distinguish the information that explains a numerical gain from the information that
supports an accuracy guarantee. Geometric-mean centering makes an extrema-based block
statistic identically one, reducing a block-relative coupling rule to Euclidean normalization;
the idealized coupling problem has the closed-form optimal factor $(1+\rho^2)^{-1/2}$.
We construct residual-budget contracts in application units and characterize exactly the row
factors compatible with a budget and coefficient limits, including incompatibility conditions,
constrained logarithmic minimax centering, integer-rounding restrictions and exact binary64
export with power-of-two factors. Budgets retain their meaning under consistent row re-emission.
Kernel-matched controls then distinguish ownership metadata
from the scaling mechanism. On 306 generated dispatch models, fixed-basis condition numbers
fall by six to eight orders of magnitude in class medians, but solver-style equilibration
recovers most of that gain and can reverse the kernel ordering. Reconstructed dispatch runs
use 30--32\% less solver work on the instances every policy completes, while a timeout reverses
the failure-sensitive ordering; the gain reproduces with a second seed but not a second solver.
In an exactly enumerated stress test, plain centering loses verified optima on all three solvers,
whereas adding a residual budget to the same rule loses none relative to that rule.
Equivalent formulations can therefore behave differently in finite precision, and external
scaling requires both mechanism attribution and an accuracy contract in original units.
